Теперь можно рассмотреть второе доказательство основной теоремы алгебры:

\begin{proof}
Пусть \(P(z)\)~--- многочлен \(n \ge 1\). Тогда он имеет хотя бы один корень в \(\Co\).
Доказательство этого факта будем вести от противного: если корней нет, то \(f(z) = 
\dfrac{1}{P(z)}\)~--- целая функция в \(\Co\), при этом \(f(z) \sim \dfrac{1}{z^n}\). \(\lim\limits_
{z \to \infty} f(z) = 0\), тогда \(\exists R_0\colon |f(z)| \le 1 \, \forall z \in \mathring{B}_
R(\infty)\). Но тогда из теоремы Лиувилля следует, что \(f(z) \equiv \mathrm{const}\), тогда и
\(P(z) \equiv \mathrm{const}\), но \(P(z)\)~--- многочлен \(n \ge 1\). Получили противоречие.

Тогда поочерёдным поиском корней \(z_k\) и делением многочлена на \((z - z_k)\) можно найти \(n\) 
корней исходного многочлена \(P(z)\).
\end{proof}


\Section{Теорема об обратной функции}

\begin{theorem} \label{29:ma}
\textbf{(МА).} Пусть \(u = u(x, y)\), \(v = v(x, y) \in C^1(G)\), где \(G \subset \R^2\)~---
область. Пусть \((x_0, y_0) \in G\) и \(\left. \begin{vmatrix} u_x & u_y \\ v_x & v_y \end{vmatrix} 
\right|_{(x_0, y_0)} \neq 0\). Обозначим \(u_0 = u(x_0, y_0)\), \(v_0 = v(x_0, y_0)\). Тогда
существуют \(B_\delta (x_0, y_0)\) и \(B_\epsilon (u_0, v_0)\), такие что:
\begin{enumerate}[wide=20pt, itemsep=0ex]
\item \(\begin{vmatrix} u_x & u_y \\ v_x & v_y \end{vmatrix} \ne 0\) в \(B_\delta (x_0, y_0)\); 
\item \(\forall (\tilde{u}, \tilde{v}) \in B_\epsilon (u_0, v_0)\) система \(\begin{cases} u(x, y) =
    \tilde{u}, \\ v(x, y) = \tilde{v} \end{cases}\) имеет в \(B_\delta (x_0, y_0)\) единственное
    решение, то есть в \(B_\epsilon (u_0, v_0)\) задана векторная функция \(x = x(u, v), y = y(u, v)
    \);
\item \(x(u, v), y(u, v) \in C^1(B_\epsilon(u_0, v_0))\).
\end{enumerate}
\end{theorem}

\begin{statement}
\textbf{(Теорема об обратной функции, локальный вариант).} Пусть:
\begin{enumerate}[wide=20pt, itemsep=0ex]
\item \(G\)~--- область в \(\Co\);
\item \(f(z)\) регулярна в \(G\);
\item \(z_0 \in G\colon f'(z_0) \ne 0\), также обозначим \(w_0 = f(z_0)\).
\end{enumerate}
Тогда существуют \(B_\delta(z_0)\) и \(B_\epsilon(w_0)\), такие что:
\begin{enumerate}[wide=20pt, itemsep=0ex]
\item \(f'(z_0) \ne 0\) в \(B_\delta (z_0)\);
\item \(\forall \hat{w} \in B_\epsilon(w_0)\) уравнение \(f(z) = \hat{w}\) имеет единственное
    решение в \(B_\delta(z_0)\), тем самым в \(B_\delta(w_0)\) задана функция \(z = z(w)\) обратная
    к \(f(z)\), то есть \(f(g(w)) \equiv w\) в \(B_\epsilon(w_0)\)
\item Эта \(g(w)\) регулярна в \(B_\epsilon(w_0)\), при этом \(g'(w) = \dfrac{1}{f'(g(w))}\).
\end{enumerate}
\end{statement}

\begin{proof}
По условию \(f(z) \in D^\infty(G)\). Пусть \(f(z) = u(x, y) + iv(x, y)\), где \(u, v \in 
C^\infty(G)\), \(z_0 = x_0 + iy_0\), \(w_0 = u_0 + iv_0\). При этом согласно утверждению 
\ref{9:Koshi_Riman} получим, что \(\begin{vmatrix} u_x & u_y \\ v_x & v_y \end{vmatrix} = u_x^2 + 
v_x^2 = |f'(z)|^2 \ne 0\) при подстановке в \(z_0\). Тогда по теореме \ref{29:ma} существуют
функции \(x = x(u, v), y = y(u, v)\) в \(B_\epsilon(w_0)\), обратные к \(u(x, y), v(x,y)\). Тогда
имеем как минимум непрерывную функцию \(g(w) = x(u, v) + iy(u, v)\). Пусть \(\hat{w} \in 
B_\epsilon(w_0)\). Докажем существование \(g'(\hat{w})\). Рассмотрим последовательность Гейне
\(w_n\) (\(w_n \to \hat{w}, w_n \ne \hat{w}\)). Пусть \(\hat{z} = g(\hat{w})\), \(z_n = g(w_n)\).
В силу непрерывности \(g(w)\) \(z_n \to \hat{z}\), а в силу взаимной однозначности области значений
\(g(w)\) и \(B_\epsilon(w_0)\) имеем: \(\lim\limits_{n \to \infty} \dfrac{g(w_n) - g(\hat{w})}{w_n -
\hat{w}} = \lim\limits_{n \to \infty} \dfrac{z_n - \hat{z}}{w_n - \hat{w}} = \dfrac{1}{\lim\limits_
{n \to \infty} \frac{w_n - \hat{w}}{z_n - \hat{z}}} = \dfrac{1}{\lim\limits_
{n \to \infty} \frac{f(z_n) - f(\hat{z})}{z_n - \hat{z}}} = \dfrac{1}{f'(\hat{z})} = \dfrac{1}
{f'(g(\hat{w}))}\).
\end{proof}


\Section{Теорема о локальной структуре отображения}

\begin{statement}
\textbf{(Теорема о локальной структуре отображения).} Пусть:
\begin{enumerate}[wide=20pt, itemsep=0ex]
\item \(f(z)\) регулярна в области \(G \subset \Co\), \(z_0 \in G\), \(w_0 = f(z_0)\);
\item \(\exists n \ge 2\colon f^{(k)}(z_0) = 0,\, \forall k \in \overline{1, n-1}, f^{(n)}(z_0) \ne
0\).
\end{enumerate}
Тогда \(\exists B_\delta(z_0), B_\epsilon(w_0)\) такие, что для любого \(\hat{w} \in \mathring{B}_
\epsilon(w_0)\) уравнение \(f(z) = \hat{w}\) имеет в \(B_\delta(z_0)\) ровно \(n\) простых корней.
\end{statement}

\begin{proof}
\begin{enumerate}[wide=20pt, itemsep=0ex]
\item Рассмотрим случай \(w_0 = 0\). Тогда \(z_0\)~--- нуль порядка \(n\) для \(f(z)\) и нуль
    порядка \(n-1\) для \(f'(z)\). Нули конечного порядка изолированы, а значит \(\exists 
    \overline{B_\delta(z_0)} \subset G\) в которой \(f(z), f'(z) \ne 0\) за исключением \(z = z_0\).
    Обозначим \(\epsilon = \min\limits_{|z - z_0| = \delta} |f(z)| > 0\). Покажем, что заключение
    теоремы выполняется для выбранных \(\delta, \epsilon\).

    Зафиксируем \(\hat{w} \in B_\epsilon (0)\). Рассмотрим в \(\overline{B_\delta (z_0)}\) функцию
    \(f(z) - \hat{w}\). Применим теорему Руше, полагая \(g(z) = - \hat{w}\). Тогда для любого
    \(\hat{z}\colon |\hat{z} - z_0| = \delta\) выполняется: \(|g(\hat{z}| = -|\hat{w}| < \epsilon =
    \min\limits_{|z - z_0| = \delta} |f(z)| < f(\hat{z}))\). По теореме Руше \(h(z) = f(z) - 
    \hat{w}\) имеет как и \(f(z)\) \(n\) нулей в \(B_\delta(z_0)\) с учётом их кратности. Пусть
    \(h(z)\) имеет нули \(z_1, \dots, z_m\) порядков \(n_1, \dots, n_m\) соответственно. Тогда 
    \(h(z_k) = 0\), но \(h'(z_k) = f'(z_k) \ne 0\) в силу выбора \(\delta\). Тогда \(z_k\)~---
    простые нули \(h(z)\), то есть уравнение \(f(z) - \hat{w} = 0\) имеет в \(B_\delta(z_0)\) ровно
    \(n\) простых корней.
\item Случай \(w_0 \ne 0\) будет рассмотрен на следующей лекции (13 лекция, 30.11).
\end{enumerate}
\end{proof}

\begin{note}
В теореме об обратной функции и теореме о локальной структуре отображения числа \(\delta\) и
\(\epsilon\) могут быть взяты меньше любого \(a < 0\).
\end{note}

\begin{statement}
\textbf{(Принцип сохранения области).} Пусть \(f(z) \ne \mathrm{const}\) регулярна в области \(D\).
Тогда \(G = f(D)\)~--- область.
\end{statement}

\begin{proof}
\begin{enumerate}[wide=20pt, label=\Roman*., itemsep=0ex]
\item Докажем открытость. Пусть \(w_0 \in G\). Тогда \(\exists z_0 \in D\colon f(z_0) = w_0\).
    \begin{enumerate}[wide=20pt, itemsep=0ex]
    \item Пусть \(f'(z_0) \ne 0\). Тогда \(\exists B_\delta(z_0) \in D, B_\epsilon(w_0) \in G\)
        такие, что для любого \(w \in B_\epsilon(w_0)\) сущесвтует и единственен \(z \in
        B_\delta(z_0)\) при котором \(f(z) = w\). Тогда \(B_\epsilon(w_0) \subset f(B_\delta(z_0))
        \subset f(D) = G\).
    \item Пусть \(\exists n \ge 2\colon f^{(k)}(z_0) = 0,\, \forall k \in \overline{1, n-1}, 
        f^{(n)}(z_0) \ne 0\). По теореме о локальной структуре отображения \(\exists B_\delta(z_0),
        B_\epsilon(w_0)\) такие, что для любого \(\hat{w} \in \mathring{B}_\epsilon(w_0)\)
        существует \(n\) решений уравнения \(f(z) = \hat{w}\) в \(B_\delta(z_0)\). Для \(w_0\) в
        \(B_\delta(z_0)\) сущетсвует по крайней мере одно решение \(z_0\) уравнения \(f(z) = w_0\),
        то есть \(B_\epsilon(w_0) \subset f(B_\delta(z_0)) \subset f(D) = G\).
    \item Пусть \(\forall n \in \N\) \(f^{(n)}(z) = 0\). Тогда \(f(z) \equiv f(z_0)\) в 
        \(B_\delta(z_0)\) и по теореме единственности получаем, что \(f(z) \equiv \mathrm{const}\)
        на всей \(D\).
    \end{enumerate}
\item Докажем линейную связность. Пусть \(w_1, w_2 \in G\). Тогда \(\exists z_1, z_2 \in D \colon
    f(z_1) = w_1, f(z_2) = w_2\). В силу связности \(D\) существует ломаная \(\gamma \subset D\),
    соединяющая \(z_1\) и \(z_2\). Обозначим \(\Gamma = f(\gamma)\). \(\Gamma\)~--- непрерывная (на
    самом деле даже КНД) кривая, соединяющая \(w_1, w_2\), лежащая в \(f(D)\). Для любого \(w \in
    \Gamma \exists B(w) \in G\), то есть существует конечное покрытие компакта \(\Gamma\). Тогда
    по лемме Гейне-Борреля выберем конечное подпокрытие из окружностей с центрами \(w^1, \dots, 
    w^n\). Соединив звеньями ломаной точки \(w_1, w^1, \dots, w^n, w_2\) в указанном порядке,
    завершаем доказательство.
\end{enumerate}
\end{proof}

\begin{definition}
Отображение \(f(z)\) называется однолистным на множестве \(E\), если из \(z_1 \ne z_2\), где 
\(z_1, z_2 \in E\), следует, что \(f(z_1) \ne f(z_2)\).
\end{definition}

\begin{example}
\(f(z) = e^z\) однолистна в \(\{|z|<3\}\), но не однолистна в \(\{|z| < 4\}\).
\end{example}

\begin{statement}
\textbf{(Теорема об обратной функции, глобальный вариант).} Пусть \(f(z)\) регулярна и однолистна в
области \(D\). Тогда:
\begin{enumerate}[wide=20pt, itemsep=0ex]
\item \(G = f(D)\)~--- область;
\item \(f'(z) \ne 0 \) для любого \(z \in D\);
\item Обратная функция \(z = g(w)\) регулярна в \(G\) и \(g'(w) = \dfrac{1}{f'(g(w))}\)
\end{enumerate}
\end{statement}

\begin{proof}
\begin{enumerate}[wide=20pt, itemsep=0ex]
\item Вытекает из принципа сохранения области, так как \(f(z) \ne \mathrm{const}\) в следствие
    однолистности.
\item От противного: пусть \(\exists z_0 \in D\colon f'(z_0) = 0\). Тогда \(\exists n \ge 2\colon 
    f^{(k)}(z_0) = 0,\, \forall k \in \overline{1, n-1}, f^{(n)}(z_0) \ne 0\). По теореме о
    локальной структуре отображения \(\exists B_\epsilon(w_0)\subset f(D)\colon w_0 = f(z_0)\) и \(\exists 
    B_\delta(z_0) \subset D\colon \forall w \in \mathring{B}_\epsilon(w_0)\) уравнение \(f(z) = w\)
    имеет \(n \ge 1\) решений. Получаем противоречие однолистности.
\item Доказательнство пункта будет изложено на следующей лекции (13 лекция 30.11).
\end{enumerate}
\end{proof}